% answers_short.tex —— 2011 年·上海卷·文科 · 简短答案速查
% 由 tools/make_short_answers.py 自 answers.tex 的 \gksolhead/\gkanswer 提取，
% 并与 2011/2011_数学真题_上海文_答案要点.md 逐题比对后生成，勿手改（改源 MD 后重跑）。

\begin{gknotice}
\textbf{使用说明}：本节只给最终答案便于快速核对，完整推导见解析版。答案为编写者独立推导，\textbf{非官方标准答案}。
\end{gknotice}

\gkpart{填空题}
\begin{itemize}[leftmargin=2.6em]
  \item[\textbf{1.}] $\{x\mid x<1\}$（即 $(-\infty,1)$）。
  \item[\textbf{2.}] $-2$。
  \item[\textbf{3.}] $-\dfrac32$。
  \item[\textbf{4.}] $\sqrt5$。
  \item[\textbf{5.}] $x+2y-11=0$。
  \item[\textbf{6.}] $(-\infty,0)\cup(1,+\infty)$。
  \item[\textbf{7.}] $3\pi$。
  \item[\textbf{8.}] $\sqrt6$。
  \item[\textbf{9.}] $\dfrac52$。
  \item[\textbf{10.}] $2$。
  \item[\textbf{11.}] $6$。
  \item[\textbf{12.}] $\dfrac{15}{2}$。
  \item[\textbf{13.}] $0.985$。
  \item[\textbf{14.}] $[-2,7]$。
\end{itemize}

\gkpart{选择题}
\begin{center}\begin{tabular}{ccccc}
\toprule
题号 & 15 & 16 & 17 & 18 \\
\midrule
答案 & A & D & A & B \\
\bottomrule
\end{tabular}\end{center}

\gkpart{解答题}
\begin{itemize}[leftmargin=2.6em]
  \item[\textbf{19.}] $z_2=4+2\mathrm i$。
  \item[\textbf{20.}] （1）$\dfrac{\sqrt{10}}{10}$；（2）$\dfrac23$。
  \item[\textbf{21.}] （1）$a,b$ 同为正时 $f(x)$ 在 $\mathbf R$ 上单调递增，同为负时在 $\mathbf R$ 上单调递减；（2）$a>0,b<0$ 时 $x<\log_{2/3}\left(-\dfrac{2b}{a}\right)$，$a<0,b>0$ 时 $x>\log_{2/3}\left(-\dfrac{2b}{a}\right)$。
  \item[\textbf{22.}] （1）$(\pm\sqrt3,0)$；（2）最大值 $5$，最小值 $\dfrac{\sqrt2}{2}$；（3）$1<m\leqslant1+\sqrt2$。
  \item[\textbf{23.}] （1）$9,15,21$；（2）$10$ 项；（3）$S_{4n}=12n^2+33n$。
\end{itemize}
